% Źródła: Eves s. 286, 290--291; Greenberg 2010, s. 206, 209; Wikipedia: Adrien-Marie Legendre.
% https://en.wikipedia.org/wiki/Adrien-Marie_Legendre

\section{Adrien-Marie Legendre} % lang-pl
\section{Adrien-Marie Legendre} % lang-it
\label{sekcja_legendre}

Trzecią wielką próbą będzie praca francuskiego analityka, która dotrze do tysięcy uczniów. % lang-pl
W tej sekcji opowiemy o Legendrze i jego podręczniku, o dwóch twierdzeniach, które faktycznie udowodni, i o trzecim, którego nie zdoła. % lang-pl
Il terzo grande tentativo sarà il lavoro di un analista francese, che raggiungerà migliaia di studenti. % lang-it
In questa sezione parleremo di Legendre e del suo manuale, dei due teoremi che dimostrerà davvero e di un terzo che non riuscirà a dimostrare. % lang-it

Adrien-Marie Legendre urodzi się 18 września 1752 roku w Paryżu i tam umrze 9 stycznia 1833 roku. % lang-pl
Jego \emph{Éléments de géométrie} (1794) przez około sto lat będą wiodącym podręcznikiem elementarnej geometrii: uprości i przestawi twierdzenia Euklidesa, a dwanaście wydań ukaże się do 1823 roku. % lang-pl
Pierwszy angielski przekład wykona w 1819 roku John Farrar z Harvardu, drugi w 1824 roku Thomas Carlyle (33 amerykańskie wydania), i tak podręcznik stanie się wzorem dla amerykańskich książek do geometrii \cite[s. 286]{eves1_1972}. % lang-pl
Adrien-Marie Legendre nascerà il 18 settembre 1752 a Parigi e vi morirà il 9 gennaio 1833. % lang-it
I suoi \emph{Éléments de géométrie} (1794) saranno per circa un secolo il testo elementare di riferimento: semplificheranno e riordineranno le proposizioni di Euclide, e dodici edizioni usciranno fino al 1823. % lang-it
La prima traduzione inglese la farà nel 1819 John Farrar di Harvard, la seconda nel 1824 Thomas Carlyle (33 edizioni americane), e così il manuale diventerà il modello dei libri americani di geometria \cite[p. 286]{eves1_1972}. % lang-it
\index[persons]{Legendre, Adrien-Marie}%
\index[persons]{Farrar, John}%
\index[persons]{Carlyle, Thomas}%

Legendre zacznie od nowa: rozważy trzy hipotezy zależnie od tego, czy suma kątów trójkąta jest mniejsza od, równa albo większa od dwóch kątów prostych. % lang-pl
Przy cichym założeniu nieskończoności prostej wyeliminuje trzecią, ale mimo powtarzanych prób nie poradzi sobie z pierwszą; kolejne próby ukażą się w kolejnych wydaniach podręcznika, a ostatnią opublikuje w 1833 roku, w roku śmierci \cite[s. 286]{eves1_1972}. % lang-pl
Legendre ricomincerà daccapo: considererà tre ipotesi a seconda che la somma degli angoli di un triangolo sia minore, uguale o maggiore di due angoli retti. % lang-it
Assumendo tacitamente l'infinità della retta eliminerà la terza, ma nonostante ripetuti tentativi non riuscirà a sbarazzarsi della prima; i tentativi successivi compariranno nelle edizioni successive del manuale, e l'ultimo lo pubblicherà nel 1833, anno della morte \cite[p. 286]{eves1_1972}. % lang-it

\begin{theorem}
[pierwsze twierdzenie Legendre'a] % lang-pl
[primo teorema di Legendre] % lang-it
\index{twierdzenie!Legendre'a (o sumie kątów)} % lang-pl
\index{teorema!di Legendre (sulla somma degli angoli)} % lang-it
	Suma kątów trójkąta nie jest większa od dwóch kątów prostych. % lang-pl
	La somma degli angoli di un triangolo non è maggiore di due angoli retti. % lang-it
\end{theorem}

Dowód (przy założeniu nieskończoności prostej i aksjomatu Archimedesa): Eves \cite[s. 290, zad. 11]{eves1_1972}. % lang-pl
Dimostrazione (assumendo l'infinità della retta e l'assioma di Archimede): Eves \cite[p. 290, es. 11]{eves1_1972}. % lang-it
Greenberg \cite[s. 206, 209]{greenberg_2010} przypisuje ten wynik Saccheriemu i Legendre'owi (zwykle nazywa się go twierdzeniem Saccheriego--Legendre'a) i dodaje, że wystarcza słabszy od aksjomatu Archimedesa aksjomat Arystotelesa (zobacz sekcję \ref{sekcja_dobre_stare}); wtedy wyklucza się suma kątów większa od $\pi$. % lang-pl
Greenberg \cite[p. 206, 209]{greenberg_2010} attribuisce questo risultato a Saccheri e a Legendre (di solito è detto teorema di Saccheri--Legendre) e aggiunge che basta l'assioma di Aristotele, più debole di quello di Archimede (si veda la sezione \ref{sekcja_dobre_stare}); allora si esclude la somma degli angoli maggiore di $\pi$. % lang-it

\begin{theorem}
[drugie twierdzenie Legendre'a] % lang-pl
[secondo teorema di Legendre] % lang-it
	Jeśli istnieje jeden trójkąt o sumie kątów równej dwóm kątom prostym, to suma kątów każdego trójkąta jest równa dwóm kątom prostym. % lang-pl
	Se esiste un triangolo con somma degli angoli uguale a due angoli retti, allora la somma degli angoli di ogni triangolo è uguale a due angoli retti. % lang-it
\end{theorem}

Dowód: Eves \cite[s. 290--291, zad. 13]{eves1_1972}; tam też wersja dla sumy mniejszej niż $\pi$. % lang-pl
Dimostrazione: Eves \cite[p. 290--291, es. 13]{eves1_1972}; ivi anche la versione per la somma minore di $\pi$. % lang-it

Trzeciego twierdzenia Legendre nie udowodni, bo ono nie jest prawdziwe: nie da się wykluczyć sumy kątów mniejszej niż $\pi$. % lang-pl
W jednym z podejść zechce zbudować trójkąt zawierający dany trójkąt co najmniej dwa razy i w tym celu założy, że przez punkt wewnątrz kąta mniejszego niż $60^\circ$ zawsze da się poprowadzić prostą przecinającą oba ramiona -- czyli coś równoważnego piątemu postulatowi (Eves \cite[s. 290, zad. 12]{eves1_1972}). % lang-pl
Il terzo teorema Legendre non lo dimostrerà, perché non è vero: non si può escludere una somma degli angoli minore di $\pi$. % lang-it
In uno dei suoi tentativi vorrà costruire un triangolo contenente un triangolo dato almeno due volte e, a tal fine, assumerà che per un punto interno a un angolo minore di $60^\circ$ si possa sempre tracciare una retta che interseca entrambi i lati -- cioè qualcosa di equivalente al quinto postulato (Eves \cite[p. 290, es. 12]{eves1_1972}). % lang-it

Jego proste i przejrzyste dowody będą tak rozpowszechnione i on sam tak znany, że w całej Europie wzbudzi zainteresowanie problemem równoległych. % lang-pl
Rekord wytrwałości w próbach dowodu postulatu należy prawdopodobnie do niego \cite[s. 286]{eves1_1972}. % lang-pl
Nic dziwnego, że sprzeczności nie znajdzie: geometria z hipotezy kąta ostrego jest równie niesprzeczna jak euklidesowa \cite[s. 286]{eves1_1972}. % lang-pl
Le sue dimostrazioni semplici e limpide, diffuse dal manuale, e la sua alta reputazione susciteranno un vivo interesse per il problema delle parallele in tutta Europa. % lang-it
Il record di perseveranza nel tentare di dimostrare il postulato spetta probabilmente a lui \cite[p. 286]{eves1_1972}. % lang-it
Non c'è da stupirsi che non trovi alcuna contraddizione: la geometria dell'ipotesi dell'angolo acuto è coerente quanto quella euclidea \cite[p. 286]{eves1_1972}. % lang-it
